% 图契约：平方反比场、立体角和点电荷奇点的关系；几何、方向和公式与正文同一约定。
\documentclass[tikz,border=5pt]{standalone}
\usepackage{fontspec,xeCJK,amsmath,amssymb}
\IfFontExistsTF{Times New Roman}{\setmainfont{Times New Roman}}{\setmainfont{TeX Gyre Termes}}
\IfFontExistsTF{Heiti SC}{\setCJKmainfont[AutoFakeBold=1.5]{Heiti SC}}{\setCJKmainfont[AutoFakeBold=1.5]{FandolHei-Regular}}
\usetikzlibrary{arrows.meta,calc,patterns}
\definecolor{ink}{HTML}{243745}
\definecolor{blue}{HTML}{275DAB}
\definecolor{green}{HTML}{177C63}
\definecolor{red}{HTML}{BD4935}
\definecolor{orange}{HTML}{AA6C16}
\definecolor{purple}{HTML}{715393}
\begin{document}
\begin{tikzpicture}[x=1cm,y=1cm,font=\fontsize{9.5}{12}\selectfont,
v/.style={-{Latex[length=2mm,width=1.3mm]},line width=.85pt},
axis/.style={-{Latex[length=1.7mm,width=1.1mm]},ink!65,line width=.55pt},
aux/.style={ink!45,dashed,line width=.55pt},
lab/.style={fill=white,inner sep=1.8pt},
card/.style={draw=ink!25,fill=blue!3,rounded corners=2pt,text width=6.1cm,minimum height=.85cm,inner sep=6pt,align=center}]
\path[use as bounding box](-8,-3.6)rectangle(8,3.65);
\node[font=\bfseries\fontsize{10.5}{13}\selectfont,text width=15.1cm,align=center]at(0.00000,3.22000){M7-H14　平方反比场、立体角和点电荷奇点的关系};
\begin{scope}[shift={(-4,0)}]
\node[text=red]at(-2.00000,0.00000){$+q$};
\draw[blue](-1.80000,0.00000)--(2.00000,-1.29968);
\draw[blue](-1.80000,0.00000)--(2.00000,1.29968);
\draw[green,line width=2pt](2,-1.3)--(2,1.3);
\node[]at(2.40000,1.60000){$dS$};
\node[]at(-0.95000,1.00000){$d\Omega$};
\draw[v,blue](-0.50000,0.00000)--(1.30000,0.00000);
\node[]at(0.40000,-0.40000){$r$};
\end{scope}
\node[card]at(4.00000,1.25000){$\mathbf E=q\mathbf e_r/(4\pi\varepsilon_0r^2)$};
\node[card]at(4.00000,-1.05000){$d\Phi_E=q\,d\Omega/(4\pi\varepsilon_0)$};
\end{tikzpicture}\end{document}
